\begin{exercises}
  \item[]\textit{这里总结我们对
    $2$-\hbox{}与 $3$-置换所用的记号。
    \begin{center}
      \begin{tabular}[t]{c|cc}
        $i$          &$1$      &$2$    \\
        \hline
        $\phi_1(i)$  &$1$      &$2$     \\
        $\phi_2(i)$  &$2$      &$1$
      \end{tabular}
      \qquad
      \begin{tabular}[t]{c|ccc}
        $i$          &$1$     &$2$   &$3$    \\
        \hline
        $\phi_1(i)$  &$1$     &$2$   &$3$    \\
        $\phi_2(i)$  &$1$     &$3$   &$2$    \\
        $\phi_3(i)$  &$2$     &$1$   &$3$    \\
        $\phi_4(i)$  &$2$     &$3$   &$1$    \\
        $\phi_5(i)$  &$3$     &$1$   &$2$    \\
        $\phi_6(i)$  &$3$     &$2$   &$1$
      \end{tabular}
    \end{center}  }
  \recommended \item 对该矩阵，求与每个 $3$-置换相联系的项。
    \begin{equation*}
      M=\begin{mat}
        1 &2 &3 \\
        4 &5 &6 \\
        7 &8 &9
      \end{mat}
    \end{equation*}
    也就是说，填写下表中的其余部分。
      \begin{center} \small
        \begin{tabular}{r|cccccc}
          置换 $\phi_i$
            &$\phi_1$ &$\phi_2$ &$\phi_3$ &$\phi_4$ &$\phi_5$ &$\phi_6$  \\
          \hline
          项 $m_{1,\phi_i(1)}m_{2,\phi_i(2)}m_{3,\phi_i(3)}$
           &$1\cdot 5\cdot 9$
           &
           &
           &
           &
           &
        \end{tabular}
      \end{center}
    \begin{answer}
      记该矩阵为~$M$。
      \begin{center}
        \begin{tabular}{r|cccccc}
          置换
            &$\phi_1$ &$\phi_2$ &$\phi_3$ &$\phi_4$ &$\phi_5$ &$\phi_6$  \\
          \hline
          项
           &$1\cdot 5\cdot 9$
           &$1\cdot 6\cdot 8$
           &$2\cdot 4\cdot 9$
           &$2\cdot 6\cdot 7$
           &$3\cdot 4\cdot 8$
           &$3\cdot 5\cdot 7$
        \end{tabular}
      \end{center}
    \end{answer}
  \recommended \item 对每个 $3$-置换~$\phi$，求 $|P_\phi|$。
    \begin{answer}
      我们可以把每个 $P_{\phi}$ 经过行对换变为单位矩阵。
      如果对换的次数是偶数，那么它的行列式为~$+1$；
      而如果对换的次数是奇数，那么行列式为~$-1$。
      \begin{center}
        \begin{tabular}{r|cccccc}
          置换
            &$\phi_1$ &$\phi_2$ &$\phi_3$ &$\phi_4$ &$\phi_5$ &$\phi_6$  \\
          \hline
          对换次数 &$0$
                  &$1$
                  &$1$
                  &$2$
                  &$2$
                  &$1$  \\
          $|P_{\phi_i}|$
           &$+1$
           &$-1$
           &$-1$
           &$+1$
           &$+1$
           &$-1$
        \end{tabular}
      \end{center}
      \textit{注。}
      在上表中，
      我们只是不断地对换，直到找到一个把我们带回到`$1,2,3$'的数为止。
      但不同的对换次数是否可能？
      如果一个人找到 $2$ 次对换而另一个人找到
      $4$ 次，这没有问题，因为两者给出的行列式都是~$+1$。
      但如果我们能用两种不同的方式对换，其中一种次数为偶数而另一种为奇数，
      那就有问题了。
      下一节将表明偶数与奇数混合的情形不可能出现。
    \end{answer}
  \item 由 $\nbyn{2}$ 公式知该行列式为~$7$。
    请用置换展开计算它。
    \begin{equation*}
      \begin{vmatrix}
        2 &3 \\
        1 &5
      \end{vmatrix}
    \end{equation*}
    \begin{answer}
      \begin{align*}
        \begin{vmat}
          2  &3 \\
          1  &5
        \end{vmat}
        &=
        2\cdot 5\cdot\deter{P_{\phi_1}}
        +
        1\cdot 3\cdot\deter{P_{\phi_2}}      \\
        &=
        2\cdot 5\cdot\begin{vmat}[r]
                           1  &0 \\
                           0  &1
                         \end{vmat}
        +
        1\cdot 3\cdot\begin{vmat}[r]
                           0  &1 \\
                           1  &0
                         \end{vmat}               \\
        &=2\cdot 5\cdot 1+1\cdot 3\cdot (-1)=7
        \end{align*}
    \end{answer}
  \item 该行列式为 $0$，因为前两行之和等于第三行。
    请用置换展开计算该行列式。
    \begin{equation*}
      \begin{vmatrix}
        -1 &0 &1 \\
        3  &1 &4 \\
        2  &1 &5
      \end{vmatrix}
    \end{equation*}
    \begin{answer}
          \begin{align*}
            \begin{vmat}[r]
              -1  &0  &1  \\
               3  &1  &4  \\
               2  &1  &5
            \end{vmat}
            &=
            \begin{aligned}[t]
              &(-1)(1)(5)\deter{P_{\phi_1}}
              +(-1)(4)(1)\deter{P_{\phi_2}}
              +(0)(3)(5)\deter{P_{\phi_3}}  \\
              &\hbox{}\quad\hbox{}
              +(0)(4)(2)\deter{P_{\phi_4}}
              +(1)(3)(1)\deter{P_{\phi_5}}
              +(1)(1)(2)\deter{P_{\phi_6}}
            \end{aligned}                        \\
            &=
            \begin{aligned}[t]
              &(-1)(1)(5)\cdot 1
              +(-1)(4)(1)\cdot (-1)
              +(0)(3)(5)\cdot (-1)  \\
              &\hbox{}\quad\hbox{}
              +(0)(4)(2)\cdot 1
              +(1)(3)(1)\cdot 1
              +(1)(1)(2)\cdot (-1)
            \end{aligned}                        \\
            &=-5+4+0+0+3-2=0
          \end{align*}
    \end{answer}
  \recommended \item
    用置换展开计算行列式。
    \begin{exparts*}
      \partsitem
        $\begin{vmat}[r]
          1  &2  &3  \\
          4  &5  &6  \\
          7  &8  &9
        \end{vmat}$
      \partsitem
        $\begin{vmat}[r]
          2  &2  &1  \\
          3  &-1 &0  \\
          -2 &0  &5
        \end{vmat}$
    \end{exparts*}
    \begin{answer}
      \begin{exparts}
        \partsitem 该矩阵是奇异的。
          \begin{align*}
            \begin{vmat}[r]
              1  &2  &3  \\
              4  &5  &6  \\
              7  &8  &9
            \end{vmat}
            &=
            \begin{aligned}[t]
              &(1)(5)(9)\deter{P_{\phi_1}}
              +(1)(6)(8)\deter{P_{\phi_2}}
              +(2)(4)(9)\deter{P_{\phi_3}}  \\
              &\hbox{}\quad\hbox{}
              +(2)(6)(7)\deter{P_{\phi_4}}
              +(3)(4)(8)\deter{P_{\phi_5}}
              +(7)(5)(3)\deter{P_{\phi_6}}
            \end{aligned}                        \\
            &=0
          \end{align*}
        \partsitem 该矩阵是非奇异的。
          \begin{align*}
            \begin{vmat}[r]
              2  &2  &1  \\
              3  &-1 &0  \\
              -2 &0  &5
            \end{vmat}
            &=
            \begin{aligned}[t]
              &(2)(-1)(5)\deter{P_{\phi_1}}
              +(2)(0)(0)\deter{P_{\phi_2}}
              +(2)(3)(5)\deter{P_{\phi_3}} \\
              &\hbox{}\quad\hbox{}
              +(2)(0)(-2)\deter{P_{\phi_4}}
              +(1)(3)(0)\deter{P_{\phi_5}}
              +(-2)(-1)(1)\deter{P_{\phi_6}}
            \end{aligned}                      \\
            &=-42
        \end{align*}
      \end{exparts}
    \end{answer}
  \recommended \item
    分别用高斯消元法与置换展开公式计算下面各行列式。
    \begin{exparts*}
      \partsitem \( \begin{vmat}[r]
                 2  &1  \\
                 3  &1
               \end{vmat}   \)
      \partsitem \( \begin{vmat}[r]
                  0  &1  &4  \\
                  0  &2  &3  \\
                  1  &5  &1
               \end{vmat}  \)
    \end{exparts*}
    \begin{answer}
      \begin{exparts}
        \partsitem 高斯消元法给出
          \begin{equation*}
             \begin{vmat}[r]
                 2  &1  \\
                 3  &1
             \end{vmat}
             =
             \begin{vmat}[r]
                 2  &1  \\
                 0  &-1/2
             \end{vmat}
             =-1
          \end{equation*}
          而置换展开给出
          \begin{equation*}
             \begin{vmat}[r]
                 2  &1  \\
                 3  &1
             \end{vmat}
             =
             \begin{vmat}[r]
                 2  &0  \\
                 0  &1
             \end{vmat} +
             \begin{vmat}[r]
                 0  &1  \\
                 3  &0
             \end{vmat}
             =
             (2)(1)\begin{vmat}[r]
                 1  &0  \\
                 0  &1
             \end{vmat} +
             (1)(3)\begin{vmat}[r]
                 0  &1  \\
                 1  &0
             \end{vmat}
             =
             -1
          \end{equation*}
        \partsitem 高斯消元法给出
          \begin{equation*}
            \begin{vmat}[r]
              0  &1  &4  \\
              0  &2  &3  \\
              1  &5  &1
            \end{vmat}
            =
            -\begin{vmat}[r]
              1  &5  &1  \\
              0  &2  &3  \\
              0  &1  &4
            \end{vmat}
            =
            -\begin{vmat}[r]
              1  &5  &1  \\
              0  &2  &3  \\
              0  &0  &5/2
            \end{vmat}
            =-5
          \end{equation*}
          而置换展开给出
          \begin{align*}
            \begin{vmat}[r]
              0  &1  &4  \\
              0  &2  &3  \\
              1  &5  &1
            \end{vmat}
            &=
            \begin{aligned}[t]
            &(0)(2)(1)\deter{P_{\phi_1}}
            +(0)(3)(5)\deter{P_{\phi_2}}
            +(1)(0)(1)\deter{P_{\phi_3}}  \\
            &\hbox{}\quad\hbox{}
            +(1)(3)(1)\deter{P_{\phi_4}}
            +(4)(0)(5)\deter{P_{\phi_5}}
            +(4)(2)(1)\deter{P_{\phi_6}}
            \end{aligned}                         \\
            &=-5
          \end{align*}
      \end{exparts}
    \end{answer}
  \recommended \item
    用置换展开公式推导 \( \nbyn{3} \) 行列式的公式。
    \begin{answer}
        仿照 \nearbyexample{ex:SamplePermExp} 可得
          \begin{align*}
             \begin{vmat}
                t_{1,1}  &t_{1,2}  &t_{1,3}  \\
                t_{2,1}  &t_{2,2}  &t_{2,3}  \\
                t_{3,1}  &t_{3,2}  &t_{3,3}
             \end{vmat}
             &=\begin{aligned}[t]
                 &t_{1,1}t_{2,2}t_{3,3}\deter{P_{\phi_1}}
                   +t_{1,1}t_{2,3}t_{3,2}\deter{P_{\phi_2}} \\
                 &\hbox{}\quad\hbox{}
                  +t_{1,2}t_{2,1}t_{3,3}\deter{P_{\phi_3}}
                             +t_{1,2}t_{2,3}t_{3,1}\deter{P_{\phi_4}} \\
                 &\hbox{}\quad\hbox{}
                  +t_{1,3}t_{2,1}t_{3,2}\deter{P_{\phi_5}}
                             +t_{1,3}t_{2,2}t_{3,1}\deter{P_{\phi_6}}
               \end{aligned}                                               \\
             &=\begin{aligned}[t]
                 & t_{1,1}t_{2,2}t_{3,3}(+1)
                   +t_{1,1}t_{2,3}t_{3,2}(-1)  \\
                 &\hbox{}\quad\hbox{}
                   +t_{1,2}t_{2,1}t_{3,3}(-1)
                    +t_{1,2}t_{2,3}t_{3,1}(+1) \\
                 &\hbox{}\quad\hbox{}
                   +t_{1,3}t_{2,1}t_{3,2}(+1)
                         +t_{1,3}t_{2,2}t_{3,1}(-1)
               \end{aligned}
          \end{align*}
     \end{answer}
  \item
    列出全部 $4$-置换。
    \begin{answer}
      下面是所有满足 $\phi(1)=1$ 的置换
      \begin{align*}
        &\phi_1=\sequence{1,2,3,4}
        \quad
        \phi_2=\sequence{1,2,4,3}
        \quad
        \phi_3=\sequence{1,3,2,4}  \\
        &\quad
        \phi_4=\sequence{1,3,4,2}
        \quad
        \phi_5=\sequence{1,4,2,3}
        \quad
        \phi_6=\sequence{1,4,3,2}
      \end{align*}
      下面是满足 $\phi(1)=1$ 的那些
      \begin{align*}
        &\phi_7=\sequence{2,1,3,4}
        \quad
        \phi_8=\sequence{2,1,4,3}
        \quad
        \phi_9=\sequence{2,3,1,4}    \\
        &\quad
        \phi_{10}=\sequence{2,3,4,1}
        \quad
        \phi_{11}=\sequence{2,4,1,3}
        \quad
        \phi_{12}=\sequence{2,4,3,1}
      \end{align*}
      下面是满足 $\phi(1)=3$ 的那些
      \begin{align*}
        &\phi_{13}=\sequence{3,1,2,4}
        \quad
        \phi_{14}=\sequence{3,1,4,2}
        \quad
        \phi_{15}=\sequence{3,2,1,4}  \\
        &\quad
        \phi_{16}=\sequence{3,2,4,1}
        \quad
        \phi_{17}=\sequence{3,4,1,2}
        \quad
        \phi_{18}=\sequence{3,4,2,1}
      \end{align*}
      以及满足 $\phi(1)=4$ 的那些。
      \begin{align*}
        &\phi_{19}=\sequence{4,1,2,3}
        \quad
        \phi_{20}=\sequence{4,1,3,2}
        \quad
        \phi_{21}=\sequence{4,2,1,3}  \\
        &\quad
        \phi_{22}=\sequence{4,2,3,1}
        \quad
        \phi_{23}=\sequence{4,3,1,2}
        \quad
        \phi_{24}=\sequence{4,3,2,1}
      \end{align*}
    \end{answer}
  \item
    置换可以看作从集合
    $\set{1,..,n}$ 到其自身的一一对应（既是单射又是满射）的函数。
    因此，每个置换都有逆。
    \begin{exparts}
      \partsitem 求每个 $2$-置换的逆。
      \partsitem 求每个 $3$-置换的逆。
    \end{exparts}
    \begin{answer}
      下列每一个都容易验证。
      \begin{exparts}
        \partsitem
          \begin{tabular}[t]{r|cc}
            置换 &$\phi_1$  &$\phi_2$ \\
             \hline
            逆     &$\phi_1$  &$\phi_2$
          \end{tabular}
        \partsitem
          \begin{tabular}[t]{r|cccccc}
            置换
              &$\phi_1$ &$\phi_2$ &$\phi_3$ &$\phi_4$ &$\phi_5$ &$\phi_6$ \\
            \hline
            逆
              &$\phi_1$ &$\phi_2$ &$\phi_3$ &$\phi_5$ &$\phi_4$ &$\phi_6$
          \end{tabular}
      \end{exparts}
    \end{answer}
  \item
     证明：\( f \) 是多重线性的当且仅当对所有
     \( \vec{v},\vec{w}\in V \) 与 \( k_1,k_2\in\Re \)，下式成立。
      \begin{equation*}
         f(\vec{\rho}_1,\dots,k_1\vec{v}_1+k_2\vec{v}_2,
           \dots,\vec{\rho}_n)
         =
         k_1f(\vec{\rho}_1,\dots,\vec{v}_1,\dots,\vec{\rho}_n)+
         k_2f(\vec{\rho}_1,\dots,\vec{v}_2,\dots,\vec{\rho}_n)
       \end{equation*}
       \begin{answer}
         对于「如果」部分，取 $k_1=k_2=1$ 即得
         \nearbydefinition{def:multilinear} 的第一个条件，
         取 $k_2=0$ 即得第二个条件。

         「仅当」部分也是常规推导。
         由
         $
           f(\vec{\rho}_1,\dots,k_1\vec{v}_1+k_2\vec{v}_2,
             \dots,\vec{\rho}_n)
         $
         出发，\nearbydefinition{def:multilinear} 的第一个条件给出
         $
           =
           f(\vec{\rho}_1,\dots,k_1\vec{v}_1,\dots,\vec{\rho}_n)+
           f(\vec{\rho}_1,\dots,k_2\vec{v}_2,\dots,\vec{\rho}_n)
         $
         再对第二个条件应用两次，即得结论。
       \end{answer}
  \item
    如果把定义中的性质~(4) 改成 \( \deter{I}=2 \)，行列式会如何变化？
    \begin{answer}
       它们都会变成原来的两倍。
    \end{answer}
  \item
    验证 \nearbycorollary{cor:ColSwapChgSign} 中的第二、第三个论断。
    \begin{answer}
      对于第二个论断，
      给定一个矩阵，先转置，再对换行，然后转置回来。
      所得结果是对换了列，而行列式改变了因子
      \( -1 \)。
      第三个论断类似：~给定
      一个矩阵，先转置，对如今已是行的东西应用多重线性，
      然后把所得的各个矩阵转置回来。
    \end{answer}
  \recommended \item
    证明：若 \( \nbyn{n} \) 矩阵的行列式非零，
    则我们能将任一列向量
    \( \vec{v}\in\Re^n \) 表示为该矩阵各列的线性组合。
    \begin{answer}
      行列式非零的 \( \nbyn{n} \) 矩阵的秩为
      \( n \)，所以它的各列构成 \( \Re^n \) 的一组基。
    \end{answer}
  \item
    \cite{Strang}
    判断真假：元素只取 $0$ 或 $1$ 的矩阵，其行列式必等于零、一或负一。
    \begin{answer}
      假。
      \begin{equation*}
        \begin{vmat}[r]
          0 &1  &1  \\
          1 &0  &1 \\
          1  &1 &0
        \end{vmat}
        =2
      \end{equation*}
     \end{answer}
  \item
    \begin{exparts}
      \partsitem 证明 \( \nbyn{5} \) 矩阵的置换
         展开公式中有 $120$ 项。
      \partsitem
         若 \( 1,2 \) 元为零，确定必为零的项有多少个？
    \end{exparts}
    \begin{answer}
      \begin{exparts}
        \partsitem 第一行元素的列下标有
          五种选择。
          于是第二行元素的列下标有
          四种选择。
          依此类推，我们得到 $5\cdot 4\cdot 3\cdot 2\cdot 1=120$。
          （另见下一题。）
        \partsitem 一旦选定了第一行的第二列，
          其余元素可以有 \( 4\cdot 3\cdot 2\cdot 1=24 \)
          种选法。
      \end{exparts}
    \end{answer}
  \item
    共有多少个 \( n \)-置换？
    \begin{answer}
       \( n\cdot(n-1)\cdots 2\cdot 1=n! \)
    \end{answer}
  \item 证明置换矩阵的逆是它的转置。
    \begin{answer}
      \cite{SchmidtSO}
      我们将证明 $P\trans{P}=I$；$\trans{P}P=I$ 的论证类似。
      $P\trans{P}$ 的 $i,j$ 元是形如
      $p_{i,k}q_{k,j}$ 的项之和，其中 $\trans{P}$ 的元素用 $q$ 记，
      即 $q_{k,j}=p_{j,k}$。
      于是 $P\trans{P}$ 的 $i,j$ 元是和
      \(
        \sum_{k=1}^n p_{i,k}p_{j,k}
      \)。
      但 $p_{i,k}$ 通常为 $0$，所以 $P_{i,k}P_{j,k}$ 通常为 $0$。
      $P_{i,k}$ 非零仅当它为 $1$，
      而此时再无其他 $i^\prime\neq i$ 使 $P_{i^\prime,k}$
      非零（$i$ 是第~$k$ 列中含 $1$ 的唯一一行）。
      换句话说，
      \begin{equation*}
        \sum_{k=1}^n p_{i,k}p_{j,k}=
        \begin{cases}
          1  &i=j  \\
          0  &\text{其他}
        \end{cases}
      \end{equation*}
      而这正是单位矩阵各元素的公式。
    \end{answer}
  \item
    矩阵 \( A \) 称为
    \definend{反对称的}\index{matrix!skew-symmetric}%
    \index{skew-symmetric}，若 \( \trans{A}=-A \)，
    如下矩阵。
    \begin{equation*}
      A=\begin{mat}[r]
          0  &3  \\
         -3  &0
        \end{mat}
    \end{equation*}
    证明：行列式非零的 \( \nbyn{n} \) 反对称矩阵
    仅在 \( n \) 为偶数时存在。
    \begin{answer}
      在
      \( \deter{A}=\deter{\trans{A}}=\deter{-A}=(-1)^n\deter{A} \)
      中指数 $n$ 必须为偶数。
    \end{answer}
  \recommended \item
    为保证 \( \nbyn{4} \) 矩阵的行列式为零，所需零元素的最少个数是多少？
    这些零元素应如何放置？
    \begin{answer}
      证明三个零的任何放置都不够是常规性的。
      四个零确实够了；把它们全放在同一
      行或同一列。
    \end{answer}
  \item
    设有 \( n \) 个数据点
    \( (x_1,y_1),(x_2,y_2),\dots\,,(x_n,y_n) \)，要找一条经过这些点的
    多项式 \( p(x)=a_{n-1}x^{n-1}+a_{n-2}x^{n-2}+\dots+a_1x+a_0 \)，
    那么我们可以代入这些点，得到一个 \( n \)~个方程/\( n \)~个未知数的
    线性方程组。
    该方程组的系数矩阵是
    \definend{Vandermonde 矩阵}。\index{matrix!Vandermonde}%
    \index{Vandermonde matrix}
    证明该系数矩阵的转置的行列式
    \begin{equation*}
      \begin{vmat}
         1       &1       &\ldots   &1       \\
         x_1     &x_2     &\ldots   &x_n     \\
         {x_1}^2 &{x_2}^2 &\ldots   &{x_n}^2 \\
                 &\vdots                     \\
         {x_1}^{n-1} &{x_2}^{n-1}   &\ldots   &{x_n}^{n-1}
      \end{vmat}
    \end{equation*}
    等于对所有满足 \( i<j \) 的下标 \( i,j\in\set{1,\dots,n} \)，
    形如 \( x_j-x_i \) 的项的乘积。
    （这表明，
    该行列式为零、线性方程组无解，
    当且仅当数据中的 \( x_i \) 不两两不同。）
    \begin{answer}
      $n=3$ 的情形表明了该怎么做。
      行倍加变换
      $-x_1\rho_2+\rho_3$ 与 $-x_1\rho_1+\rho_2$
      给出
      \begin{multline*}
        \begin{vmat}
          1     &1     &1     \\
          x_1   &x_2   &x_3   \\
          x_1^2 &x_2^2 &x_3^2
        \end{vmat}
        =
        \begin{vmat}
          1     &1             &1             \\
          x_1   &x_2           &x_3           \\
          0     &(-x_1+x_2)x_2 &(-x_1+x_3)x_3
        \end{vmat}                                           \\
        =
        \begin{vmat}
          1     &1             &1             \\
          0     &-x_1+x_2      &-x_1+x_3      \\
          0     &(-x_1+x_2)x_2 &(-x_1+x_3)x_3
        \end{vmat}
      \end{multline*}
      再作行倍加变换 $x_2\rho_2+\rho_3$ 便得
      所要的结果。
      \begin{equation*}
        =
        \begin{vmat}
          1     &1             &1                     \\
          0     &-x_1+x_2      &-x_1+x_3              \\
          0     &0             &(-x_1+x_3)(-x_2+x_3)
        \end{vmat}
        =(x_2-x_1)(x_3-x_1)(x_3-x_2)
      \end{equation*}
    \end{answer}
  \item
    我们可以把矩阵分成
    \definend{块}，\index{block matrix}\index{matrix!block} %
    如此处，
    \begin{equation*}
      \begin{pmat}{rr|r}
          1  &2   &0  \\
          3  &4   &0  \\  \cline{1-3}
          0  &0   &-2
       \end{pmat}
    \end{equation*}
    其中显示出四个块：左上与右下的方形的 $\nbyn{2}$ 与 $\nbyn{1}$ 块，
    以及右上与左下的零块。
    证明：若一个矩阵可以分块为
    \begin{equation*}
      T=
      \begin{pmat}{c|c}
          J   &Z_2  \\  \cline{1-2}
          Z_1 &K
       \end{pmat}
    \end{equation*}
    其中 $J$ 与 $K$ 是方阵，而 $Z_1$ 与 $Z_2$ 全为零，
    则 \( \deter{T}=\deter{J}\cdot\deter{K} \)。
    \begin{answer}
      设 \( T \) 为 \( \nbyn{n} \) 的，
      \( J \) 为 \( \nbyn{p} \) 的，
      \( K \) 为 \( \nbyn{q} \) 的。
      应用置换展开公式
      \begin{equation*}
        \deter{T}=
        \sum_{\text{\scriptsize 置换 }\phi}
                t_{1,\phi(1)}t_{2,\phi(2)}\dots t_{n,\phi(n)}
                \deter{P_{\phi}}
      \end{equation*}
      由于 \( T \) 的右上部分全为零，如果一个
      \( \phi \) 的前
      \( p \) 个列编号 \( \phi(1),\dots,\phi(p) \) 中至少有一个取自 \( p+1,\dots,n \)，
      那么由 \( \phi \) 产生的项为 \( 0 \)
      （例如，若 \( \phi(1)=n \) 则
      \( t_{1,\phi(1)}t_{2,\phi(2)}\dots t_{n,\phi(n)} \)
      为 \( 0 \)）。
      所以上式化简为对具有两段的全部置换求和：
      先重排 \( 1,\dots,p \)，随后是
      \( p+1,\dots,p+q \)。
      为看出这给出 \( \deter{J}\cdot\deter{K}  \)，将其展开。
      \begin{multline*}
         \bigg[\sum_{\substack{\text{\scriptsize 置换 }\phi_1 \\
                              \text{\scriptsize 作用于 } 1,\dots,p}}
               t_{1,\phi_1(1)}\cdots t_{p,\phi_1(p)}
               \deter{P_{\phi_1}} \bigg]                                  \\
         \cdot
         \bigg[\sum_{\substack{\text{\scriptsize 置换 }\phi_2 \\
                              \text{\scriptsize 作用于 } p+1,\dots,p+q}}
               t_{p+1,\phi_2(p+1)}\cdots t_{p+q,\phi_2(p+q)}
               \deter{P_{\phi_2}}                          \bigg]
      \end{multline*}
    \end{answer}
  \item
    证明：对任一 \( \nbyn{n} \) 矩阵 \( T \)，使矩阵 \( T-rI \) 的
    行列式为零的相异实数 \( r \) 至多有
    \( n \) 个
    （我们将在第五章用到这一结果）。
    \begin{answer}
      $n=3$ 的情形表明了会发生什么。
      \begin{equation*}
        \deter{T-rI}
        =
        \begin{vmat}
          t_{1,1}-x  &t_{1,2}   &t_{1,3}  \\
          t_{2,1}    &t_{2,2}-x &t_{2,3}  \\
          t_{3,1}    &t_{3,2}   &t_{3,3}-x
        \end{vmat}
      \end{equation*}
      置换展开中的每一项都有三个取自
      矩阵元素的因子（例如，$(t_{1,1}-x)(t_{2,2}-x)(t_{3,3}-x)$
      与 $(t_{1,1}-x)(t_{2,3})(t_{3,2})$），所以该行列式
      可以表示为 $x$ 的 $3$ 次多项式。
      这样的多项式至多有 $3$ 个根。

      一般地，置换展开表明
      行列式是一些项的和，每一项
      有 \( n \) 个因子，给出一个 $n$ 次多项式。
      \( n \) 次多项式至多有 \( n \) 个根。
    \end{answer}
  \puzzle \item
    \cite{MathMag63Q307}
    九个正的数字可以按 \( 9! \) 种方式排成 \( \nbyn{3} \) 阵列。
    求这些阵列的行列式之和。
    \begin{answer}
      \answerasgiven
      当行列式的两行对换时，行列式的
      符号改变。
      当一个三阶行列式的行作置换时，得到 \( 3 \)
      个正的与 \( 3 \) 个绝对值相等的
      负行列式。
      于是 \( 9! \) 个行列式分成 \( 9!/6 \) 组，每一组
      之和为零。
    \end{answer}
  \item
    \cite{MathMag63Q237}
    证明
    \begin{equation*}
      \begin{vmat}
        x-2  &x-3  &x-4  \\
        x+1  &x-1  &x-3  \\
        x-4  &x-7  &x-10
      \end{vmat}=0.
    \end{equation*}
    \begin{answer}
      \answerasgiven
      当任一列的元素分别从另外两列的元素中减去后，所得
      行列式的两列元素成比例，所以行列式为零。
      也就是，
      \begin{equation*}
        \begin{vmat}
          2  &1    &x-4  \\
          4  &2    &x-3  \\
          6  &3    &x-10
        \end{vmat}=
        \begin{vmat}
          1    &x-3  &-1   \\
          2    &x-1  &-2   \\
          3    &x-7  &-3
        \end{vmat}=
        \begin{vmat}
          x-2  &-1   &-2   \\
          x+1  &-2   &-4   \\
          x-4  &-3   &-6
        \end{vmat}=0.
      \end{equation*}
    \end{answer}
  \puzzle \item
    \cite{Monthly49p33}
    设 \( S \) 为一个三阶幻方的整数元素之和，\( D \) 为把该方阵视为
    行列式时的值。
    证明 \( D/S \) 是整数。
    \begin{answer}
      \answerasgiven
      设
      \begin{equation*}
        \begin{array}{ccc}
          a  &b  &c  \\
          d  &e  &f  \\
          g  &h  &i
        \end{array}
      \end{equation*}
      其幻和为 \( N=S/3 \)。
      那么
      \begin{align*}
        N
        &=(a+e+i)+(d+e+f)+(g+e+c)   \\
        &\quad\text{}-(a+d+g)-(c+f+i)=3e
      \end{align*}
      且 \( S=9e \)。
      因此，把行与列相加，
      \begin{equation*}
        D=
        \begin{vmat}
          a  &b  &c  \\
          d  &e  &f  \\
          g  &h  &i
        \end{vmat}
        =\begin{vmat}
          a  &b  &c  \\
          d  &e  &f  \\
         3e  &3e &3e
        \end{vmat}
        =\begin{vmat}
          a  &b  &3e \\
          d  &e  &3e \\
         3e  &3e &9e
        \end{vmat}
        =\begin{vmat}
          a  &b  &e  \\
          d  &e  &e  \\
          1  &1  &1
        \end{vmat}S.
      \end{equation*}
    \end{answer}
  \puzzle \item
    \cite{Monthly31p355}
    证明 Pascal 三角形左上角 \( n^2 \) 个元素构成的行列式
    \begin{equation*}
      \begin{array}{cccccc}
        1  &1  &1  &1  &.  &.  \\
        1  &2  &3  &.  &.      \\
        1  &3  &.  &.  &   &   \\
        1  &.  &.  &   &   &   \\
        .                      \\
        .
      \end{array}
    \end{equation*}
    的值为 $1$。
    \begin{answer}
      \answerasgiven
      记所论行列式为 \( D_n \)，记
      \( i \) 行 \( j \) 列的元素为 \( a_{i,j} \)。
      那么由元素的构造法则我们有
      \begin{equation*}
        a_{i,j}=a_{i,j-1}+a_{i-1,j},
        \qquad a_{1,j}=a_{i,1}=1.
      \end{equation*}
      从最后一对行开始，将 \( D_n \) 的每一行从其后一行中减去。
      经过 \( n-1 \) 次这样的减法，上面的等式表明元素
      \( a_{i,j} \) 被元素 \( a_{i,j-1} \) 所取代，而第一列中除
      \( a_{1,1}=1 \) 外的所有元素都变成零。
      现在从最后一对列开始，将每一列从其后一列中减去。
      经过这一过程，如上面的关系式所示，元素 \( a_{i,j-1} \) 被
      \( a_{i-1,j-1} \) 所取代。
      这两个运算的结果是把 \( a_{i,j} \) 换成
      \( a_{i-1,j-1} \)，并使第一行与
      第一列中的每个元素都化为零。
      于是 \( D_n=D_{n+i} \)，从而
      \begin{equation*}
        D_n=D_{n-1}=D_{n-2}=\dots=D_2=1.
      \end{equation*}
    \end{answer}
\end{exercises}
